\chapter{News and Event Trading}\label{ch:hf:news-and-event-trading}

At half past eight on the first Friday of the month a number is released. Within a millisecond futures trade in Chicago, and the market makers
who did not want to be part of it pulled their quotes a second before. In this chapter's release model the fastest tier of traders, 50
microseconds from the number, takes 88\% of everything the release's stale quotes give away and the next tier the rest; a market maker who leaves
its quotes loses 58 tick-lots a release, and one that pulls them and comes back five seconds later earns 29\% more over the next two minutes than
one that never left.

\section{Scheduled numbers and the release race}

Economic releases (One Quant Book 2, chapter 31) arrive at known times with a consensus forecast; the price moves with the data surprise. Their
publishers control who sees the number first.

\begin{definition}[Release lock-up]\label{def:hf:news-and-event-trading:lockup}
A \emph{release lock-up}\index{release lock-up} is an arrangement in which a statistical agency or data publisher gives journalists or data
vendors the number before its public release, under embargo and without outside communication, so that stories and feeds are ready at the release
time.
\end{definition}

\begin{definition}[Release race]\label{def:hf:news-and-event-trading:race}
A \emph{release race}\index{release race} is the \omterm{def:hf:latency-arbitrage-and-its-defence:arb}{latency race} (chapter 9) triggered by a scheduled release: traders who receive and parse the
number first take the quotes that the number has made stale, in every instrument that depends on it.
\end{definition}

The race's outcome depends on who gets the number first, and on how quickly the quotes it makes stale are withdrawn. Hu, Pan and Wang studied a
period from 2007 to June 2013 when select high-speed traders received the Michigan index of consumer sentiment two seconds before its public release:
trading was highly concentrated in those two seconds and price discovery took less than 200 milliseconds. Comparing with other releases and with
the period after the arrangement ended, they concluded that the tiered release may have reduced, rather than enhanced, the informational
advantage of the fastest traders, and made price discovery more efficient.

\begin{dated}{2026-09}{Early access, a hacked headline and a tweet}\label{dat:hf:news-and-event-trading:events}
On 8 July 2013 the New York Attorney General announced that Thomson Reuters had agreed to discontinue immediately the practice of providing
high-frequency traders with certain market-moving consumer survey results two seconds before its other subscribers. On 23 April 2013, just after
1 p.m., a false tweet from the Associated Press's hacked account reported explosions at the White House; the Dow Jones Industrial Average, the
Nasdaq and the S\&P 500 spiked downward for a few minutes until traders realised it was false. On 29 September 2018 the SEC announced that Elon Musk
and Tesla would each pay a \$20 million penalty and that Musk would step down as chairman, settling charges over his tweets that he could take Tesla
private at \$420 a share with funding secured.
\end{dated}

\section{Machine-readable headlines}

Machine-readable news (One Quant Book 8, chapter 17) sends headlines and numbers in fields a program can parse without reading. The trader that
parses first wins the race; the trader that parses wrongly loses it. A strategy that captures 2 ticks when it reads a headline right and loses 10
when it reads it wrong earns 1.76 ticks a headline at a 2\% error rate and loses at 20\%: misparsing, not speed, is often the binding constraint
once a firm is fast enough.

\section{Social media and unscheduled events}

Unscheduled events arrive without a consensus, often through channels built for people rather than machines, and sometimes false. The
April 2013 hack in the dated box moved the whole US market for minutes; a tweet by a company's chief executive moved its stock and ended in a
settlement. A market maker cannot win a race it does not know has started; what it can do is recognise that one has, and protect its quotes.

\section{The market maker's event protocol}

\begin{definition}[Event protocol]\label{def:hf:news-and-event-trading:protocol}
An \emph{event protocol}\index{event protocol} is a market maker's written schedule for a known event: when to cancel or widen quotes before it,
which inputs must be updated before quoting resumes, when and how wide to re-enter, and which limits apply until normal quoting resumes.
\end{definition}

The chapter's model (\cref{lst:hf:news-and-event-trading:race}): the number's surprise moves the future's fair value by three ticks per standard
deviation; ten levels of fifty lots each, from providers who update 20 milliseconds after the release, plus our market maker's twenty lots at each of
the first three levels; four tiers of takers at 0.05, 1, 5 and 50 milliseconds, two hundred lots each. A lot taken $k$ ticks from the old price
earns $J-k$ ticks for a move of $J$.

\begin{center}
\small
\begin{tabular}{@{}lrrrr@{}}
\toprule
takers' latency & 0.05 ms & 1 ms & 5 ms & 50 ms\\
\midrule
share of the stale quotes' value captured & 88.1\% & 11.6\% & 0.3\% & 0\\
\bottomrule
\end{tabular}
\end{center}

The race is winner-take-most: 223 tick-lots of stale value per release, almost all to the first tier. Our market maker's quotes sit behind the
others' (it was not first to the level) and still lose 57.5 tick-lots a release on average (\cref{lst:hf:news-and-event-trading:mmloss}): at \$12.50 a
tick, \$719.

The protocol's second half is re-entry. After the release the flow is heavy (twenty-one times normal at first, decaying over thirty seconds) and
toxic (a mark-out of three ticks at first, decaying over five seconds). Quoting from $t$ seconds until two minutes after the release earns
(\cref{fig:hf:news-and-event-trading:reentry}): 872 tick-lots from the release itself, 911 from half a second, 1\,052 from five seconds, and less
after that as the volume fades. A market maker that stays through the release earns $872-58=814$; one that pulls and re-enters at five seconds earns
1\,052.

\begin{omfigure}
\begin{tikzpicture}
\begin{axis}[width=11cm,height=5.2cm,xlabel={\small re-entry time after the release (seconds)},ylabel={\small tick-lots until two minutes},xmin=0,
  xmax=62,ymin=0,ymax=1150,tick label style={font=\footnotesize},grid=major,grid style={gray!25},table/col sep=comma]
\addplot[omThm,thick,mark=*,mark size=1.3pt] table[x=t,y=value]{figdata/hft/18-news-and-event-trading/reentry.csv};
\draw[omDef,dashed] (axis cs:0,814) -- (axis cs:62,814);
\node[font=\scriptsize,anchor=south west,omDef] at (axis cs:30,814) {staying through the release};
\end{axis}
\end{tikzpicture}

\omcaption{Spread less mark-outs earned by quoting from the re-entry time until two minutes after a release, in tick-lots, when post-release flow
starts at twenty-one times normal and decays over thirty seconds and its mark-out starts at three ticks and decays over five; dashed, a market maker
that never pulled its quotes, after its loss in the race. Data: \texttt{hf\_news.reentry\_curve}.}\label{fig:hf:news-and-event-trading:reentry}
\end{omfigure}

\section{Risk controls for event trading}

Event trading concentrates risk in seconds. The controls of chapter 27 apply with event-specific settings: position and loss limits that tighten
before a release, a check that the parsed number is within a plausible range of the consensus before a taking strategy acts on it, a maximum
size per headline, and a kill switch that a human can reach within the event's window. The market maker's protocol is itself a risk control: an
automated cancel at a scheduled time, verified before the event, not left to a trader's reflexes.

\section{Strategy files}

\begin{strategyfile}{Macro-release race in futures}\label{strat:hf:news-and-event-trading:macro}
\sfield{Who pays you, and why} Providers whose quotes are stale for the milliseconds after a scheduled number.
\sfield{Instruments and venues} Rates, equity index and currency futures; the release's official feed or a vendor's.
\sfield{Signal} The surprise against consensus, mapped to each instrument's move from past releases.
\sfield{Sizing and execution} Take the stale levels inside the expected move, largest first; exit into the post-release flow.
\sfield{Costs} Feed and co-location at the release point and the exchanges; fees; the wrong-sign risk if the mapping is wrong.
\sfield{How it dies} The race: the first tier takes 88\% of the value in the model; the second-fastest has little.
\sfield{Horizon, capacity, infrastructure} Microseconds to seconds; capacity from stale depth, which protocols shrink.
\sfield{Backtest honestly} The release's exact timestamp at the trader's location, the depth that was resting then, and the fastest competitor.
\sfield{Sources} Hu, Pan and Wang (2017); this chapter.
\end{strategyfile}

\begin{strategyfile}{Machine-readable earnings headline race}\label{strat:hf:news-and-event-trading:earnings}
\sfield{Who pays you, and why} Providers in the stock and its options whose quotes lag the headline.
\sfield{Instruments and venues} The stock, its options, its sector fund; a machine-readable news feed.
\sfield{Signal} Parsed fields (earnings, guidance) against consensus.
\sfield{Sizing and execution} Small sizes, scaled by parsing confidence; stop on any field outside its plausible range.
\sfield{Costs} The feed; misparsing (a 20\% error rate turns 2-tick captures and 10-tick losses into a loss).
\sfield{How it dies} Parsing errors and faster competitors on the same feed.
\sfield{Horizon, capacity, infrastructure} Milliseconds to minutes.
\sfield{Backtest honestly} Feed timestamps at receipt, not at publication; the headlines as first sent, including corrections.
\sfield{Sources} One Quant Book 8, chapter 17.
\end{strategyfile}

\begin{strategyfile}{Social-media headline reaction}\label{strat:hf:news-and-event-trading:social}
\sfield{Who pays you, and why} Slower traders reacting to the same post; or, on false posts, the traders who acted on them.
\sfield{Instruments and venues} Index futures and the stocks named; social media feeds.
\sfield{Signal} Posts from verified sources, with credibility scores and confirmation from other sources.
\sfield{Sizing and execution} Tiny until confirmed; fading moves that lack confirmation, as after the April 2013 hack.
\sfield{Costs} False positives; the feed.
\sfield{How it dies} Hacked or false accounts; regulatory action against manipulative posting.
\sfield{Horizon, capacity, infrastructure} Seconds to minutes.
\sfield{Backtest honestly} Posts as they appeared, including deleted ones, with the timestamp of receipt.
\sfield{Sources} The dated box.
\end{strategyfile}

\begin{strategyfile}{Market-maker event protocol: pull, widen, re-enter}\label{strat:hf:news-and-event-trading:protocol}
\sfield{Who pays you, and why} The heavy post-release flow, once its toxicity has decayed.
\sfield{Instruments and venues} Every instrument the market maker quotes that the release moves.
\sfield{Signal} The release calendar; after the release, the decay of mark-outs.
\sfield{Sizing and execution} Cancel before the release; re-enter after the toxic seconds (five in the model) at a wider spread, narrowing as
mark-outs decay.
\sfield{Costs} The spread foregone in the seconds out of the market.
\sfield{How it dies} Unscheduled events, for which there is no calendar.
\sfield{Horizon, capacity, infrastructure} Seconds; an automated scheduler verified before each event.
\sfield{Backtest honestly} Post-release fills and mark-outs by second after the release, not averaged over the day.
\sfield{Sources} This chapter: 1\,052 tick-lots re-entering at five seconds against 814 staying.
\end{strategyfile}

\section{Tutorial: half past eight}\label{tut:hf:news-and-event-trading:tut}

\begin{tutorial}
\textbf{Goal.} Measure who captures a release's stale quotes, what a market maker loses by staying, and when it should come back.
\textbf{End state:} the table and \cref{fig:hf:news-and-event-trading:reentry}.
\begin{tutsteps}
\item \textbf{The race}: tiers in latency order take stale levels inside the move.
\omcode{code/firm/newsrace/firm_newsrace.py}{32}{52}{Fastest first, each tier up to its capacity; a lot $k$ ticks from the old price earns $J-k$.}\label{lst:hf:news-and-event-trading:race}
\item \textbf{The market maker's loss} from the back of each level's queue.
\omcode{code/firm/newsrace/firm_newsrace.py}{55}{62}{Our lots are taken only after the others' at each level.}\label{lst:hf:news-and-event-trading:mmloss}
\item \textbf{Twenty thousand releases} with standard normal surprises (\texttt{hf\_news.base}).
\item \textbf{Re-entry}: spread less decaying mark-outs from each re-entry time (\texttt{hf\_news.reentry\_curve}).
\end{tutsteps}
\textbf{What to change next.} Replay the release on \texttt{firm.exchsim} with the providers as agents that react after their own latencies;
add a second instrument that the number moves; let the surprise's mapping be estimated with error.
\end{tutorial}

\section{Build: the release race module}\label{bld:hf:news-and-event-trading:newsrace}

\begin{build}
\bfield{Purpose} Simulate scheduled releases, the race for stale quotes, and a market maker's \omterm{def:hf:news-and-event-trading:protocol}{event protocol}.
\bfield{Interface} \texttt{move(z, beta)}, \texttt{race(J, stale\_lots, tiers, stale\_ms)}, \texttt{mm\_loss(J, mm\_lots, other\_lots, tiers,
stale\_ms)}, \texttt{reentry(t)}, \texttt{simulate(n, seed)}, \texttt{misparse(p\_wrong, capture\_right, loss\_wrong)}. Unscheduled news:
\texttt{firm.newsevent} (Book 8).
\bfield{Rules} Ticks and lots; takers arriving after the providers' reaction get nothing; our lots are behind the others'.
\bfield{Acceptance tests} \texttt{code/firm/newsrace/tests/}: a 4-tick race by hand, including a tier too slow to take anything; the market
maker's loss by hand on both sides; re-entry after the toxic seconds beats re-entry at once; shares sum to one and the slowest tier gets
nothing; the misparsing value by hand.
\bfield{Stretch} Providers as agents on \texttt{firm.exchsim}; several instruments; a learned post-release mark-out curve.
\end{build}

\begin{omsources}
\item G.~X.~Hu, J.~Pan, J.~Wang, Early peek advantage? Efficient price discovery with tiered information disclosure, \emph{Journal of Financial
Economics} 126(2), 2017, 399--421.
\item New York State Attorney General, press release on Thomson Reuters, 8 July 2013.
\item US Securities and Exchange Commission, press release 2018-226, 29 September 2018.
\item NBC News, report on the Associated Press Twitter account hijack, 23 April 2013.
\end{omsources}

\section{Exercises}

\begin{exercise}[$\star$]\label{exo:hf:news-and-event-trading:1}
A release moves the future four ticks. A taker buys 10 lots at each of the first three levels above the old price. What does it earn?
\end{exercise}

\begin{exercise}[$\star$]\label{exo:hf:news-and-event-trading:2}
A headline strategy captures 2 ticks when right and loses 10 when wrong. At what error rate does it break even?
\end{exercise}

\begin{exercise}[$\star$]\label{exo:hf:news-and-event-trading:3}
At \$12.50 a tick, what is 57.5 tick-lots, and what is the protocol's advantage of 238 tick-lots a release worth over twelve monthly releases?
\end{exercise}

\begin{exercise}[$\star\star$]\label{exo:hf:news-and-event-trading:4}
Why does the fastest tier take almost everything?
\end{exercise}

\begin{exercise}[$\star\star$]\label{exo:hf:news-and-event-trading:5}
Why is re-entering at once worse than re-entering after five seconds, and why is re-entering after thirty seconds worse again?
\end{exercise}

\begin{exercise}[$\star\star$]\label{exo:hf:news-and-event-trading:6}
Hu, Pan and Wang found that early access for a few may have improved price discovery. Give the argument.
\end{exercise}

\begin{exercise}[$\star\star\star$]\label{exo:hf:news-and-event-trading:7}
\emph{Coding.} Rerun \texttt{firm.newsrace.simulate} with the providers updating after 0.5 ms instead of 20. What happens to the tiers' shares and to
the value captured per release?
\end{exercise}

\begin{exercise}[$\star\star\star$]\label{exo:hf:news-and-event-trading:8}
\emph{Find the flaw.} ``We backtested our release strategy on the vendor's timestamps and it captures 3 ticks on every release.''
\end{exercise}

\section{Problem: Half Past Eight}

\begin{problem}[{Weekend problem --- half past eight}]\label{pb:hf:news-and-event-trading:1}
A futures market maker and several takers face a monthly release.

\medskip\noindent\textbf{Part I --- Releases.}
\begin{enumerate}
\item Define a \omterm{def:hf:news-and-event-trading:lockup}{release lock-up} and a \omterm{def:hf:news-and-event-trading:race}{release race}.
\item What did Hu, Pan and Wang find?
\item Summarise the dated box's three events.
\item Why does misparsing matter more than speed for many firms?
\end{enumerate}

\medskip\noindent\textbf{Part II --- The race.}
\begin{enumerate}[resume]
\item Describe the model: move, stale levels, tiers.
\item Give the tiers' shares of the stale value.
\item Why is our market maker's loss computed from the back of the queue?
\item What does it lose a release, in tick-lots and dollars?
\end{enumerate}

\medskip\noindent\textbf{Part III --- The protocol.}
\begin{enumerate}[resume]
\item Define an \omterm{def:hf:news-and-event-trading:protocol}{event protocol}.
\item Give the re-entry curve's values at 0, 0.5, 5 and 30 seconds.
\item Compare staying with pulling and re-entering at five seconds.
\item Which risk controls belong in the protocol?
\end{enumerate}

\medskip\noindent\textbf{Part IV --- The verdict.}
\begin{enumerate}[resume]
\item State the \emph{named result}: the share of the release's move captured by each latency tier, and the market maker's loss with and without
its \omterm{def:hf:news-and-event-trading:protocol}{event protocol}.
\item Who pays for the race?
\item What would faster providers change?
\item How should the protocol handle unscheduled events?
\item Why must a backtest use the receipt timestamp?
\item Which strategy file is most exposed to false information?
\item How would you estimate the post-release mark-out curve?
\item In one sentence: what does a market maker's \omterm{def:hf:news-and-event-trading:protocol}{event protocol} buy?
\end{enumerate}
\end{problem}

\section{Interview questions}

\begin{interviewq}[$\star$ \iqroles{trader}]\label{iq:hf:news-and-event-trading:1}
Payrolls are in ten minutes. What do you do with your quotes, and when do you come back?
\end{interviewq}

\begin{interviewq}[$\star\star$ \iqroles{developer}]\label{iq:hf:news-and-event-trading:2}
Design the path from a machine-readable release to an order, and list where it can fail.
\end{interviewq}

\begin{interviewq}[$\star\star$ \iqroles{researcher}]\label{iq:hf:news-and-event-trading:3}
How would you map a release's surprise to the move of ten different instruments, and test the mapping?
\end{interviewq}

\begin{interviewq}[$\star\star$ \iqroles{risk}]\label{iq:hf:news-and-event-trading:4}
A headline strategy made money on 95\% of events and lost ten times its average gain on the rest. How do you set its limits?
\end{interviewq}

\begin{interviewq}[$\star\star$ \iqroles{trader}]\label{iq:hf:news-and-event-trading:5}
A tweet says a company has received a takeover offer. What do you check before trading, and in what order?
\end{interviewq}

\begin{interviewq}[$\star\star\star$ \iqroles{researcher}]\label{iq:hf:news-and-event-trading:6}
With flow decaying as $e^{-s/T_1}$ and mark-outs as $e^{-s/T_2}$, find the re-entry time that maximises the market maker's expected earnings.
\end{interviewq}
